\section{Contar los tipos pagando los radios}
Un tipo de $v$ sitios etiquetados prescribe cuatro particiones de sus etiquetas,
una por familia. Sus igualdades son ecuaciones lineales en $2v$ coordenadas.
Sean $h$ el número de índices RC distintos y $D$ la dimensión de todas las
igualdades. Sólo contamos tipos realizables. Escribimos
\[
 d=h-D,\quad b=p-D-1,\quad L_s=(2v-D)-(v-c).
\]
El rango archivado da $b\ge0$. Un bosque de igualdades sobre sitios da $v-c$
ecuaciones independientes por eliminación de hojas. Se tiene $0\le L_s\le b+2$,
pues $p=v+c+1,v+c,v+c-1$. Todo sitio tiene vecino diagonal, luego $d\ge1$.
Una base de las igualdades diagonales tras identificar RC cubre los $h$ índices,
a razón de cuatro por ecuación: $h\le4d$. Así $d+b-1\ge(v-4)/3$.
\paragraph{Cuenta ponderada.}
Para $v,c,L_s$ fijos,
\[
 \frac1{v!}\sum_{\text{tipos etiquetados}}\prod_i d_i
       \le \frac{C^v(4v^2)^{L_s}}{c!}.
\]
Para verlo se asigna cada sitio a una fibra donde alcanza su máximo $d_i$,
rompiendo empates en un orden fijo. En la fibra $f$ escogida, $a_f\ge1$ sitios
asignados tienen $d_i=D_f=|f\cap U|$. Se unen sus $a_f$ miembros a un nodo auxiliar
de color $f$ (se guarda la familia, no la coordenada). Las estrellas son disjuntas.
Se añaden igualdades entre sitios, sin cerrar ciclos, hasta conectar cada componente
original. Esto es posible porque todas las incidencias que conectan componentes
parciales se pueden realizar con un par de sitios de la misma fibra.
Las igualdades internas de cada estrella pueden representarse por un árbol entre
sus sitios. Tras esa sustitución, el bosque tiene $v-c$ ecuaciones independientes.
Se eligen $L_s$ igualdades adicionales hasta completar una base del espacio de
ecuaciones. Sus dos extremos y su familia cuestan a lo sumo $(4v^2)^{L_s}$.
El código recupera las cuatro particiones por pertenencia al espacio generado:
una igualdad de tipo pertenece al espacio, y una igualdad allí contenida se cumple
en la realización y por tanto está en el tipo. Esto impide colisiones entre tipos.
Se enraíza cada componente en un sitio y se juntan las $c$ raíces a una superraíz.
Quedan a lo sumo $2v+1$ nodos, un alfabeto fijo de colores de aristas y nodos,
y $v$ etiquetas distintas de sitio. Hay a lo sumo $C^v v!$ árboles planos de esta
forma, antes de las igualdades adicionales. Se puede fijar un código por tipo.
Ordenar los hijos produce códigos planos distintos: los sitios están etiquetados y
cada nodo auxiliar se reconoce por sus miembros asignados, no vacíos y disjuntos.
La superraíz aporta $c!$ órdenes y cada estrella al menos $(a_f-1)!$.
La razón entre el peso del tipo y este número de ordenaciones está acotada por
\[
 \frac1{c!}\prod_f\frac{D_f^{a_f}}{(a_f-1)!}
 \le \frac{2^v}{c!}\exp\left(\sum_fD_f\right)\le \frac{2^ve^{4v}}{c!}.
\]
Se usó $D^a/a!\le e^D$, $a\le2^a$ y que cada sitio participa en cuatro fibras.
La división por $v!$ en la suma final corresponde a las ordenaciones de cada soporte
real, incluido cuando distintas ordenaciones presentan el mismo tipo.
